% !PRINTED 295
\weberpage{295}
%<<CONT
\gesperrt{this equation can become reducible by successive adjunction of roots
of cyclic equations only if it is imprimitive.}

We can state something further about the manner of the reduction, its possibility
being presupposed. If $P_\nu$ is the first group of the sequence (8) that is a
normal divisor of $P$, then by Theorem III, §.\,7, just cited, with the correct
arrangement of the composition series:
\weq{9}{j_\lambda = j_{\lambda+1} \cdots = j_{\nu+1},}
hence all equal to the same prime, and the resolvent $\chi = 0$, by which the
group $P$ is reduced to $P_\nu$, consequently has a metacyclic group.

If we denote by $A, B, \ldots, S$ the systems of intransitivity of the group
$P_\nu$, whose number shall be $s$, and set
\weq{10}{m = r\, s,}
then by adjunction of a root of $\chi = 0$ the function $f(x)$ will split into
$s$ factors of degree $r$.

The $A, B, \ldots, S$ are systems of imprimitivity of $P$ (Bd. I, §.\,158, 2.).

If we denote by $Q$ the totality of all permutations of $P$ which leave the
individual systems $A, B, \ldots, S$ in their places and interchange only the
elements of the systems among themselves, then $Q$, as we saw in §.\,158 of the
first volume, is a normal divisor of $P$, and by adjunction of the roots of an
auxiliary equation $\varphi(y) = 0$ of degree $s$, $f(x)$ splits into $s$
factors of degree $r$, to which the individual systems $A, B, \ldots, S$ belong
as roots:
\[
f(x) = f(x, y_\alpha)\, f(x, y_\beta)\, f(x, y_\gamma) \ldots
\]
Since by the auxiliary equation $\varphi(y) = 0$ the group $P$ is reduced to
$Q$, the group of $\varphi(y) = 0$ is isomorphic with $P/Q$. On the other hand,
however, the intransitive group $P_\nu$ is also a divisor of $Q$, since the
systems $A, B, \ldots, S$ are not displaced by $P_\nu$, and we can therefore
apply the theorem §.\,73, 3. to the three groups $P, Q, P_\nu$. Since the index
series of $P/P_\nu$ by hypothesis consists entirely of primes, by that
theorem the same holds of $P/Q$. This group is metacyclic, and hence the
auxiliary equation $\varphi(y) = 0$ is also metacyclic. There follows, then,
the theorem:
%CONT>>
